\scanpage{58}
\mspage{17}
\[
3R=\{\deleted{\unsure{\text{avec}}}\deleted{\unsure{\varnothing}}\ x'-x\mid x,x'\in\Delta'\}.
\]
Mais \deleted{soit $\widetilde\Delta$ notons que $\widetilde\Delta'$} notons que
\[
-\Delta'=\{-(3\delta\deleted{\unsure{+}}-\lambda)=\lambda-3\delta\mid\delta\in\Delta\}
\]
est aussi invariant par $W$ et contenu dans $S_{-12}$, d'ailleurs \deleted{\unsure{tel}} $\Delta'\cap(-\Delta')=\varnothing$ car $3\delta-\lambda=\lambda-3\delta'$ implique $3(\delta+\delta')=2\lambda=2(3\eta-\zeta)$ or $2\lambda=6\eta-2\zeta$ n'est pas divisible par $3$ dans $\widehat E$ !

\deleted{Notons que $W$} \unsure{Ceci|Cela} montre que $-\operatorname{id}_E\notin W$, et comme c'est $\in\widetilde W$, on trouve
\[
\widetilde W=W\times\underbrace{\mathbb Z/2\mathbb Z}_{\substack{\text{engendré par}\\\operatorname{id}_E}}
\]
\deleted{\unsure{$\Delta$}} La réunion $\widetilde\Delta^{\unsure{+|=}}=\Delta'\cup(-\Delta')$ est donc \deleted{\unsure{intrinsèqut|intrinsèque}} \inserted{canoniqt} associé au syst. de racines (indépendamment de son obtention par un $\Delta$) et il se décompose en deux orbites sous $W$. Le choix d'une orbite équivaut au choix du syst. de racines. On aimerait l'expliciter, \inserted{ainsi que $\widetilde\Delta$,} via la géométrie interne du syst. de racines.
% archive sheet number: 57
